% ==================================================================================================================================
% Introduction à l'étude Quantitative

En pratique, on ne peut que très rarement résoudre formellement une équation différentielle. 
On va donc cherche, à partir de la fonction qui la définit, de déduire des propriétés sur les solutions. 

% ==================================================================================================================================
% Introduction à l'étude Quantitative

\section{Équations Autonomes}

\begin{definition}[Équation Autonome]
    Une équation différentielle est dite autonomme si elle est de la forme : 
        \[ (A) : y'(t) = f(y(t)) \] 
    Autrement dit si la variable temporelle (i.e $t$) n'apparaît pas dans $f$. 
\end{definition}

On peut interpréter une équation autonome en physique comme une équation à la contrainte de temps n'apparaît pas. 

Posons maintenant quelques définitions supplémentaires... 

\begin{definition}[Champ de vecteurs]
    Un champ de vecteurs est une fonction continue $f : \Omega \subset \R^n \longrightarrow \R^n$ où $\Omega$ est un ouvert. 
    Intuitivement, elle associe à chaque point de l'espace $\Omega$ un vecteur. 
\end{definition}

\begin{definition}[Courbe Intégrale]
    Une courbe intégrale d'un champ de vecteur $f : \Omega \longrightarrow \R^n$ est une \textbf{courbe paramétrée} 
    $y : I \subset \R \longrightarrow \Omega$ telle que $y$ est solution maximale de l'équation autonome :
        \[ (E) : y'(t) = f(y(t)) \quad \forall i \in I \] 
\end{definition}

\begin{definition}[Trajectoire]
    Une trajectoire de $f : \Omega \longrightarrow \R^n$ est l'image d'une courbe intégrale de $f$. 
\end{definition}

\begin{definition}[Point Stationnaire]
    On dit que $ x_0 \in \Omega$ est un point stationnaire ou point d'quilibre si $f(x_0) = 0$. 
    Un point stationnaire est un point où le champ de vecteur est nul. 
\end{definition}

Ainsi, un champ de vecteur va permettre de modéliser un système en mouvement tel que l'écoulement d'un fluide. 
Une trajectoire de ce champ de vecteur sera donc la trajectoire d'une particule dans le fluide. 

\begin{proposition}
    Soit $y$ une courbe intégrale d'une équation autonome $(A)$ alors si $c \in \R$ alors la courbe :
        \[ y_c : 
            \begin{cases}
                I + c \longrightarrow \R^n \\ 
                t \longmapsto y(t - c)
            \end{cases} \] 
    est une courbe intégrale (i.e solution maximale de $A$). Elle a la même trajectoire que $y$ à translation près. 
\end{proposition}

\begin{proposition}
    Soit $(A)$ une équation autonome définie par une fonction $f$. Supposons que $f$ est localement lispchitzienne. 
    Soient :
        \begin{align*}
            y_1 : I_1 \longrightarrow \R^n \\ 
            y_2 : I_2 \longrightarrow \R^n 
        \end{align*}
    deux solutions maximales de $(A)$ telles que 
        \[ \exists t_1, t_2 \in I_1, I_2, \quad y_1(t_1) = y_2(t_2) \] 
    \underline{alors} $y_1$ et $y_2$ sont égales à translation près. 
\end{proposition}

\begin{definition}[Orbite]
    Toutes les courbes intégrales $y$ d'une équation autonome $f$ vérifiant le même problème de Cauchy ont la même trajectoire. 
    Cette trajectoire est appelée \textbf{orbite}. 
    On appelle \textbf{portrait de phase} la partition de $\Omega$ en orbites.  
\end{definition}

\begin{prop}[Orbites]
    Deux orbites sont soient disjointes soient égales. 
\end{prop}

